%=====================================================================
\chapter{Changing a Schema That Is Already Full}\label{ch:evolution}
%=====================================================================
\locator{3}{Part III --- From Diagram to Schema}

\begin{outcomes}
By the end of this chapter you will be able to:
\begin{tight}
  \item \textbf{Classify} a proposed schema change as cheap, expensive or
        impossible, and \textbf{say} what decides which.
  \item \textbf{Use} the four forms of \sqlkw{ALTER TABLE} SQLite supports, and
        \textbf{state} the limits of each.
  \item \textbf{Perform} the twelve-step rebuild for a change
        \sqlkw{ALTER TABLE} cannot make.
  \item \textbf{Write} a migration as a numbered, repeatable script, and
        \textbf{record} the schema version in the database itself.
  \item \textbf{Backfill} a new \sqlkw{NOT NULL} column on a table that already
        has rows.
\end{tight}
\end{outcomes}

\begin{prereq}
\chref{ddl}: \sqlkw{CREATE TABLE} and \sqlkw{STRICT}. \chref{mapping} for the
schema you are about to change --- this chapter modifies the
\texttt{university.db} you built there, and produces
\texttt{v2\_constraints.sql}.
\end{prereq}

\begin{keyterms}
schema evolution $\bullet$ migration $\bullet$ \sqlkw{ALTER TABLE} $\bullet$
\sqlkw{ADD COLUMN} $\bullet$ \sqlkw{RENAME TO} $\bullet$ \sqlkw{RENAME COLUMN}
$\bullet$ \sqlkw{DROP COLUMN} $\bullet$ the twelve-step rebuild $\bullet$
backfill $\bullet$ \sqlkw{PRAGMA} \texttt{user\_version} $\bullet$
forward migration $\bullet$ rollback
\end{keyterms}

\section{The Chapter the Outline Does Not Name}

HEC's course content lists nothing about changing a schema --- \term{schema
evolution}, in the textbooks' phrase. It is here anyway,
for a reason you can check against your own timetable: this course carries one
database through sixteen weeks. In week~5 it had seven tables; by
\chref{normalforms} it has more, by \chref{views} it has views and triggers,
and by \chref{indexes} it has indexes. Every one of those is a schema change
made to a database that already has rows in it.

A student who can only create a schema from scratch has to start again every
time the design improves --- and in a real system, starting again means losing
the data.

\begin{conceptbox}[Three kinds of change, and what decides the cost]
\begin{tabular}{@{}L{3.6cm}L{4.6cm}L{5.8cm}@{}}
\textbf{Kind} & \textbf{Examples} & \textbf{Cost} \\
\textbf{Cheap} & add a nullable column; rename; add an index &
  a change to the \emph{schema only}. The rows are untouched. \\
\textbf{Expensive} & change a type; add a constraint; split a table &
  every row must be read and rewritten. \\
\textbf{Impossible} & recover a fact nobody recorded &
  no amount of work helps; only a new requirement and future data. \\
\end{tabular}

\smallskip
\textbf{What decides it is whether existing rows must change.} Adding a
nullable column does not: every existing row's new value is null, and null
needs no writing. Adding \sqlkw{NOT NULL} does, because every existing row
needs a value that is not there.

\smallskip
The third row is the one worth remembering. \chref{ermodel} said a
\texttt{address} stored whole could never answer ``who lives in Bahawalpur''.
That is not expensive --- the city was never recorded separately and splitting
the column cannot invent it. Design errors in Part~II become \emph{impossible}
here, which is why Part~II took four chapters.
\end{conceptbox}

\section{What \sqlkw{ALTER TABLE} Can Do}

\begin{definitionbox}[The four forms, and the limits of each]
\begin{tight}
  \item \sqlkw{ALTER TABLE} $t$ \sqlkw{RENAME TO} $u$ --- renames the table,
        and since 3.25 updates every reference to it in other tables' foreign
        keys, views and triggers.
  \item \sqlkw{ALTER TABLE} $t$ \sqlkw{RENAME COLUMN} $a$ \sqlkw{TO} $b$ ---
        renames a column, updating references.
  \item \sqlkw{ALTER TABLE} $t$ \sqlkw{ADD COLUMN} $c$ \dots\ --- adds a
        column at the \emph{end}. It may not be \sqlkw{PRIMARY KEY} or
        \sqlkw{UNIQUE}; it may not have a non-constant default; and it may be
        \sqlkw{NOT NULL} \textbf{only if it has a default}.
  \item \sqlkw{ALTER TABLE} $t$ \sqlkw{DROP COLUMN} $c$ --- removes a column,
        provided nothing depends on it: not a primary key, not indexed, not
        mentioned in a constraint, view or trigger.
\end{tight}

\smallskip
That is the complete list. \textbf{There is no \sqlkw{ALTER COLUMN}}, so a
type cannot be changed, a constraint cannot be added, and a column cannot be
moved --- not by \sqlkw{ALTER TABLE}.
\end{definitionbox}

\begin{examplebox}[The three cheap changes, and the one that catches people]
\begin{lstlisting}[language=sql]
-- Cheap: a nullable column. Instant on any number of rows.
ALTER TABLE student ADD COLUMN email TEXT;

-- Cheap: renames.
ALTER TABLE student RENAME COLUMN cnic TO national_id;

-- Refused: NOT NULL with no default, on a table with rows.
ALTER TABLE student ADD COLUMN email TEXT NOT NULL;
--   Cannot add a NOT NULL column with default value NULL
\end{lstlisting}

\smallskip
\textbf{The third is not an arbitrary restriction.} Every existing row would
need a value for \texttt{email} and there is none, so the table would be in a
state its own schema forbids the instant the statement finished. SQLite
refuses rather than produce that.

\smallskip
The two honest ways out are both in the next section, and the dishonest one is
\sqlkw{DEFAULT} \texttt{''} --- which satisfies the constraint with a value
that means nothing, and is \chref{relational}'s magic value wearing a new hat.
\end{examplebox}

\section{The Twelve-Step Rebuild}

For everything else --- changing a type, adding a \sqlkw{CHECK}, adding
\sqlkw{STRICT}, reordering columns --- there is one procedure, and it is in
SQLite's own documentation. It is worth doing by hand once.

\begin{worked}[Adding a constraint to a table that is already full]
The requirement: \texttt{student.semester} must be between 1 and 8. The table
exists, has 4{,}000 rows, and has no such \sqlkw{CHECK}.

\medskip
\textbf{Step 0 --- find out whether the data already satisfies it.} Before
anything else:
\begin{lstlisting}[language=sql]
SELECT roll_no, semester FROM student
 WHERE semester NOT BETWEEN 1 AND 8 OR semester IS NULL;
\end{lstlisting}
If this returns rows, the migration will fail at step 5 and you need a
decision from the university first --- not a cleverer script. \textbf{This is
the step people skip, and it is the one that turns a ten-minute migration into
an afternoon.}

\medskip
\textbf{Steps 1--12, as a script:}
\begin{lstlisting}[language=sql]
PRAGMA foreign_keys = OFF;          -- 1. and NOT inside the transaction
BEGIN;                              -- 2. all of it, or none of it

CREATE TABLE student_new (          -- 3. the table you wish you had
  roll_no   TEXT    PRIMARY KEY NOT NULL,
  name      TEXT    NOT NULL,
  programme TEXT    NOT NULL,
  semester  INTEGER NOT NULL CHECK (semester BETWEEN 1 AND 8),
  cnic      TEXT    UNIQUE
) STRICT;

INSERT INTO student_new                -- 4. copy, naming every column
  SELECT roll_no, name, programme, semester, cnic FROM student;

DROP TABLE student;                    -- 5. the old one
ALTER TABLE student_new RENAME TO student;   -- 6. take its name

CREATE INDEX idx_student_name ON student(name);  -- 7. rebuild indexes
-- 8. recreate views and triggers that referenced it

PRAGMA foreign_key_check;           -- 9. did anything become an orphan?
COMMIT;                             -- 10.
PRAGMA foreign_keys = ON;           -- 11.
PRAGMA integrity_check;             -- 12.
\end{lstlisting}

\medskip
\textbf{Four things about this that are not obvious.}
\begin{tight}
  \item \textbf{\sqlkw{PRAGMA} foreign\_keys goes \emph{outside} the
        transaction}, both times. Inside one it is silently ignored --- which
        is \chref{constraints}'s pitfall arriving in a place where the
        consequence is a dropped foreign key rather than a rejected insert.
  \item \textbf{Name every column in the \sqlkw{INSERT}.} \sqlkw{SELECT *}
        depends on column order, which is exactly what you may be changing.
  \item \textbf{Indexes, views and triggers do not survive \sqlkw{DROP
        TABLE}.} Step 7 is not optional and is the most commonly forgotten.
  \item \textbf{It is $O(n)$.} Every row is read and written. On this
        database that is instant; on a real one it is the difference between a
        maintenance window and a Tuesday.
\end{tight}
\end{worked}

\dsfig{%
\begin{tikzpicture}[
  box/.style={rounded corners=3pt, draw=#1, line width=1pt, fill=#1!7,
              text=#1!60!black, font=\scriptsize\bfseries, align=center,
              text width=2.9cm, minimum height=0.8cm},
  arr/.style={-{Stealth[length=2.2mm]}, primary!55, line width=0.9pt},
  q/.style={font=\scriptsize\itshape, text=neutral, align=center},
  tag/.style={font=\scriptsize\bfseries, align=center}
]
  \node[box=primary] (start) at (0,2.6) {the change\\you want};
  \node[q, text width=4.6cm] (q1) at (0,1.2)
       {does any \emph{existing row} have to change?};
  \node[box=greenhl] (cheap) at (-4.6,-0.3) {ALTER TABLE\\{\scriptsize
       \mdseries schema only, instant}};
  \node[q, text width=4.4cm] (q2) at (2.2,-0.3)
       {is the information already in the database?};
  \node[box=goldhl] (reb) at (0.4,-2.1) {the 12-step\\rebuild};
  \node[box=accent]  (imp) at (5.0,-2.1) {\textbf{impossible}};

  \draw[arr] (start) -- (q1);
  \draw[arr] (q1) -- node[tag, above, text=greenhl!60!black, pos=0.55] {no}
             (cheap);
  \draw[arr] (q1) -- node[tag, above right, text=neutral, pos=0.5] {yes}
             (q2);
  \draw[arr] (q2) -- node[tag, left, text=goldhl!70!black] {yes} (reb);
  \draw[arr] (q2) -- node[tag, right, text=accent] {no} (imp);

  \node[q, anchor=north, text width=3.6cm, text=greenhl!60!black]
       at (-4.6,-1.1) {add a nullable column, rename, add an index};
  \node[q, anchor=north, text width=3.6cm, text=goldhl!70!black]
       at (0.4,-2.9) {change a type, add a \sqlkw{CHECK}, split a table,
        backfill};
  \node[q, anchor=north, text width=3.6cm, text=accent]
       at (5.0,-2.9) {recover a fact nobody ever recorded --- a Part~II error};
\end{tikzpicture}}%
{Two questions decide what a schema change costs. The right-hand branch is the
one Part~II existed to prevent: no procedure recovers information that was
never stored.}{evolution}

\section{Migrations as Files}

\begin{definitionbox}[A migration]
A \term{migration} is a numbered script that moves the database from one schema
version to the next. It is kept in the repository beside the code, it is run
once per database, and it is never edited after it has been run anywhere.

\smallskip
\sqlkw{PRAGMA} \texttt{user\_version} is a 32-bit integer SQLite stores in the
database file and otherwise ignores completely. It is yours to use, and it is
how a script knows which migrations a given file has already had.
\end{definitionbox}

\begin{examplebox}[The pattern, and it is four lines]
\begin{lstlisting}[language=sql]
-- v2_constraints.sql -- run only against a version 1 database
PRAGMA foreign_keys = OFF;
BEGIN;

-- ... the changes ...

PRAGMA user_version = 2;      -- the last thing before COMMIT
COMMIT;
PRAGMA foreign_keys = ON;
\end{lstlisting}

and the check before running it:
\begin{lstlisting}[language=bash]
sqlite3 university.db "PRAGMA user_version;"
# 1  -> run v2.   2 -> already applied, do nothing.
\end{lstlisting}

\smallskip
\textbf{Why the version is set inside the transaction.} If the migration
fails half-way, the rollback takes the version number with it, and the
database is honestly still at version 1. Set it afterwards and a failed
migration leaves a database that claims to be something it is not --- which is
worse than the failure.
\end{examplebox}

\begin{alertbox}[Four rules that are about people, not about SQL]
\begin{tight}
  \item \textbf{Never edit a migration that has run anywhere.} Somebody else's
        database already has the old version applied; changing the file means
        two databases with the same version number and different schemas, and
        nothing will ever tell you.
  \item \textbf{Take a backup first, and check that you can restore it.} For
        SQLite this is \texttt{copy university.db university.db.bak} and then
        actually opening the copy. An untested backup is a hope.
  \item \textbf{Make the change in two deployments when a column is going
        away.} Stop writing it, ship, confirm nothing broke, then drop it.
        Dropping and deploying together means the rollback is also a
        migration.
  \item \textbf{Write the \emph{reverse} migration, even if you never run
        it.} Writing it is how you find out that the change is irreversible
        --- and finding that out before you run it is the entire point.
  \end{tight}
\end{alertbox}

\begin{pitfall}
\begin{tight}
  \item \textbf{Adding \sqlkw{NOT NULL} without a default to a full table.}
        Refused, and correctly: every existing row would violate it.
  \item \textbf{Satisfying it with \sqlkw{DEFAULT} \texttt{''}.} A magic value
        that will be counted and compared as though it meant something.
  \item \textbf{Setting \sqlkw{PRAGMA} foreign\_keys inside the transaction.}
        Silently ignored, and in a rebuild the consequence is a dropped
        foreign key.
  \item \textbf{\sqlkw{INSERT} \dots\ \sqlkw{SELECT *} in a rebuild.} It
        depends on the column order you are changing.
  \item \textbf{Forgetting to recreate indexes, views and triggers.} They go
        with the dropped table and nothing reminds you.
  \item \textbf{Skipping step 0.} Check the data satisfies the new constraint
        \emph{before} writing the migration.
  \item \textbf{Editing a migration after it has run.} Two databases, one
        version number, different schemas.
\end{tight}
\end{pitfall}

\begin{lab}[Migrate the database you built in Chapter 9]
This lab produces \texttt{code/schema/v2\_constraints.sql}.

\begin{lstlisting}[language=bash]
# --- 1. Back it up, and prove the backup opens ----------------------
cp university.db university.db.bak
sqlite3 university.db.bak ".tables"
sqlite3 university.db "PRAGMA user_version;"     # 0 or 1
\end{lstlisting}
\begin{lstlisting}[language=sql]
-- --- 2. The cheap change: a nullable column ---------------------
ALTER TABLE student ADD COLUMN email TEXT;
SELECT roll_no, email FROM student LIMIT 3;
-- Instant, whatever the row count. Nothing was rewritten.

-- --- 3. The same column, NOT NULL -------------------------------
ALTER TABLE student ADD COLUMN email2 TEXT NOT NULL;
-- Record the error. Explain it in terms of the existing rows.

-- --- 4. Step 0 of any rebuild: does the data already comply? ----
SELECT COUNT(*) FROM student WHERE semester NOT BETWEEN 1 AND 8;
-- If this is not 0, stop and decide what the right value is.

-- --- 5. Backfill, so NOT NULL becomes possible ------------------
UPDATE student SET email = roll_no || '@iub.edu.pk' WHERE email IS NULL;
SELECT COUNT(*) FROM student WHERE email IS NULL;
-- Now every row has a value, and only now is the constraint addable.
\end{lstlisting}
\begin{lstlisting}[language=bash]
# --- 6. Write the migration -----------------------------------------
# Put steps 1-12 into code/schema/v2_constraints.sql, adding to the
# student table:  email TEXT NOT NULL, and a CHECK that semester is
# between 1 and 8. Set PRAGMA user_version = 2 before COMMIT.
\end{lstlisting}
\begin{lstlisting}[language=sql]
-- --- 7. Run it --------------------------------------------------
.read code/schema/v2_constraints.sql
PRAGMA user_version;
.schema student
PRAGMA integrity_check;
PRAGMA foreign_key_check;

-- --- 8. Prove the new constraint is live ------------------------
INSERT INTO student VALUES
  ('2021-IT-999','Test','BSIT',99,NULL,NULL,NULL,NULL,'x@y.z');
-- Must fail. Which constraint caught it?

-- --- 9. Prove the old data survived -----------------------------
SELECT COUNT(*) FROM student;
-- Compare with the count before step 6. They must match.

-- --- 10. Write the reverse migration ----------------------------
-- Write v2_rollback.sql, which undoes step 6. Do not run it.
-- Then answer: is it a true reverse? What cannot be undone?
\end{lstlisting}

\smallskip
\textbf{What to notice.} Step 10 is the exercise. The forward migration added a
\sqlkw{CHECK} and a column; the reverse can drop both --- but the
\emph{backfilled email addresses} are now indistinguishable from real ones,
and no rollback can tell them apart. \textbf{A migration that writes data is
not reversible, whatever the script says}, and knowing which of your
migrations are in that class is the whole of the skill.
\end{lab}

\begin{checkpoint}
\begin{tightnum}
  \item Classify each as cheap, expensive or impossible, with the reason:
        adding a nullable \texttt{email}; adding \texttt{email} \sqlkw{NOT
        NULL}; changing \texttt{semester} from \texttt{TEXT} to
        \texttt{INTEGER}; splitting \texttt{address} into house, street and
        city on a table where only the whole address was ever stored.
  \item A migration adds a \sqlkw{NOT NULL} column by backfilling it with a
        computed value. Explain why the rollback script is not a true reverse,
        and say what you would record at migration time to make it one.
  \item \sqlkw{PRAGMA} \texttt{user\_version} is set inside the transaction
        rather than after it. Say what goes wrong if it is set after, and
        describe the state the database would be left in.
\end{tightnum}
\end{checkpoint}

\begin{chaptersummary}
\begin{tight}
  \item A schema change is \textbf{cheap} if no existing row must change,
        \textbf{expensive} if every row must be rewritten, and
        \textbf{impossible} if it needs information nobody recorded.
  \item The third class is a Part~II design error arriving late. No procedure
        recovers a fact that was never stored.
  \item SQLite's \sqlkw{ALTER TABLE} does four things: \sqlkw{RENAME TO},
        \sqlkw{RENAME COLUMN}, \sqlkw{ADD COLUMN}, \sqlkw{DROP COLUMN}.
        \textbf{There is no \sqlkw{ALTER COLUMN}}, so types and constraints
        cannot be changed by it.
  \item \sqlkw{ADD COLUMN} may be \sqlkw{NOT NULL} only with a default,
        because every existing row needs a value --- and a default of
        \texttt{''} is a magic value, not a solution.
  \item Everything else is the \textbf{twelve-step rebuild}: new table, copy
        naming every column, drop, rename, \textbf{rebuild indexes, views and
        triggers}, check, commit.
  \item \sqlkw{PRAGMA} \texttt{foreign\_keys} goes \emph{outside} the
        transaction. Inside one it is ignored, and in a rebuild that costs you
        a foreign key.
  \item A \textbf{migration} is a numbered script, kept in the repository, run
        once, and \textbf{never edited after it has run anywhere}.
  \item \sqlkw{PRAGMA} \texttt{user\_version} records which migrations a file
        has had. Set it \emph{inside} the transaction, so a failed migration
        leaves an honest version number.
  \item Back up first and open the backup. Drop a column in a second
        deployment. Write the reverse migration even if you never run it.
  \item \textbf{A migration that writes data is not reversible}, whatever the
        rollback script claims --- backfilled values cannot be told from real
        ones afterwards.
\end{tight}
\end{chaptersummary}

\begin{reviewq}
  \item \textbf{(MCQ)} Which of these can SQLite's \sqlkw{ALTER TABLE} do?
        (a)~change a column's type; (b)~add a \sqlkw{CHECK} constraint;
        (c)~rename a column; (d)~reorder columns.
  \item \textbf{(MCQ)} \sqlkw{ADD COLUMN} \texttt{x TEXT NOT NULL} on a table
        with rows is: (a)~instant; (b)~refused unless a default is given;
        (c)~refused always; (d)~allowed but leaves nulls.
  \item \textbf{(Short)} Why must \sqlkw{PRAGMA} \texttt{foreign\_keys = OFF}
        appear before \sqlkw{BEGIN} rather than after it?
  \item \textbf{(Short)} Name two things that are lost by \sqlkw{DROP TABLE}
        during a rebuild and must be recreated explicitly.
  \item \textbf{(Short)} Give a schema change that is impossible, and say what
        decision in Part~II would have prevented it.
  \item \textbf{(Applied)} Write the full twelve-step migration that changes
        \texttt{enrolment.grade} from an unconstrained \texttt{TEXT} column to
        one with \sqlkw{CHECK} \texttt{(grade IN ('A','B','C','D','F'))}.
        Include step 0, say what you would do if step 0 returned 40 rows with
        a grade of \texttt{'I'}, and state which parts of your script are not
        reversible.
  \item \textbf{(Applied)} Your university decides, in week 12, that an
        instructor may co-teach: \ent{teaches} becomes M:N. Describe the
        complete migration --- new table, backfill from the existing
        \texttt{section.emp\_id}, what to do with the old column, and in how
        many deployments you would ship it. Then say which single step is
        irreversible and why.
\end{reviewq}
